\chapter{流形学习}
\label{chap:manifold}
\begin{quote}\small
\textit{本章摘要。}流形学习假定高维观测由少量连续因素控制。本章从线性子空间基线出发，分别推导保持测地距离、局部重构和图平滑的算法，再讨论 t-distributed Stochastic Neighbor Embedding (t-SNE)、Uniform Manifold Approximation and Projection (UMAP) 与流形正则化。低维图形只是一种表示，不能直接证明故障类别可分。
\end{quote}

上一章解决了数据分散时如何协同训练，却没有保证传递给模型的距离和表示本身合理。相同故障在不同转速下可能相距很远，不同故障也可能因传感器尺度相似而靠得很近。本章把问题从“数据在哪里”转向“数据之间什么关系应被保留”：如果观测的高维变化主要由少量连续因素驱动，能否在低维坐标中保留对判断有用的结构？

把许多传感器读数画到平面上，可以有不同的画法。PCA 尽量保留线性重构能力，Isometric Mapping (Isomap) 保留沿邻域图行走的距离，Locally Linear Embedding (LLE) 保留局部重构关系，谱方法关注图上的平滑变化；概率可视化更重视谁与谁相邻，并不保证远处两点的距离准确。本章分别推导这些目标。看图时还要问：聚在一起的是同一种故障，还是同一种转速？换一台设备后，这种表示是否仍有助于识别故障？

\section{几何假设与邻域构造}

% auxiliary-resource-guide
本章末的“辅助学习材料”列出与核心方法对应的讲解和实践入口，可在阅读相关推导后，按所列小节对照学习。

\begin{figure}[htbp]
\centering
\begin{tikzpicture}[x=0.9cm,y=0.75cm]
 \draw[gray!35] (0,0) grid (6,4);
 \draw[blue!75!black, very thick] plot[smooth] coordinates {(0.4,0.7)(1.1,1.5)(1.8,2.8)(2.7,3.1)(3.5,2.0)(4.3,0.9)(5.4,1.2)};
 \foreach \x/\y in {0.4/0.7,1.1/1.5,1.8/2.8,2.7/3.1,3.5/2.0,4.3/0.9,5.4/1.2} {\fill[red!70!black] (\x,\y) circle (2pt);}
 \draw[orange!80!black, dashed] (1.1,1.5) circle (0.55cm);
 \draw[orange!80!black, dashed] (3.5,2.0) circle (0.55cm);
 \node[font=\scriptsize, anchor=west] at (6.15,3.2) {局部邻域};
 \node[font=\scriptsize, anchor=west] at (6.15,2.65) {高维观测};
 \node[font=\scriptsize, anchor=west] at (6.15,2.1) {低维流形};
\end{tikzpicture}
\caption{流形学习把高维观测看作嵌入空间中的低维结构；邻域大小决定局部几何是否可靠。}
\label{fig:manifold-local-geometry}
\end{figure}
图\ref{fig:manifold-local-geometry}用曲线上的邻域说明局部距离与整体嵌入形状的区别；邻域过大会跨越不同局部区域，过小则可能使几何估计缺乏样本支持。
设$x=g(z)+\epsilon\in\R^d$，$z\in\R^r$且$r\ll d$，$g$在局部光滑。目标可以是恢复坐标、保持邻域或改善预测，这三者不是同一个优化问题。不同载荷、转速和传感器尺度会改变距离；先按训练数据确定标准化和度量，再构造$k$近邻或半径邻域图。

邻域太小会使图断裂，太大会跨越流形折叠形成捷径。互为近邻与并集近邻的连通性不同，应记录图的连通分量和边长分布。故障突变可能跨越不连续机制，不能因为原始欧氏距离近就强迫标签平滑。

\section{线性基线与测地距离}
\paragraph{PCA 与核 PCA。}
对中心化数据矩阵$X\in\R^{n\times d}$，PCA 求
\begin{equation}
 \min_{V^{\mathsf T}V=I}\|X-XVV^{\mathsf T}\|_F^2,
\end{equation}
其解由$X^{\mathsf T}X$的主特征向量组成。坐标为$Z=XV$，截断奇异值分解可避免显式构造大协方差矩阵。高方差方向不一定携带故障信息，低方差异常可能被截断。核 PCA 在核矩阵上中心化后求特征向量，允许非线性表示，但需要选择核与处理新样本映射\cite{scholkopf1998kpca}。

\paragraph{PCA 的重构推导与尺度选择。}
令$V\in\R^{d\times r}$列正交，$P=VV^{\mathsf T}$满足$P^2=P=P^{\mathsf T}$。因此
\begin{align}
 \|X-XP\|_F^2
 &=\operatorname{tr}[(I-P)X^{\mathsf T}X(I-P)]\\
 &=\operatorname{tr}(X^{\mathsf T}X)-\operatorname{tr}(V^{\mathsf T}X^{\mathsf T}XV).
\end{align}
最小化重构误差等价于最大化最后的迹。设$X^{\mathsf T}X=U\Lambda U^{\mathsf T}$，特征值降序排列，把$V$展开到$U$基底后，每个特征方向的被保留权重介于$0$与$1$，且总和为$r$；最大值于是为前$r$个特征值之和。若$X=U_X\Sigma V_X^{\mathsf T}$，坐标为$XV_{X,r}=U_{X,r}\Sigma_r$，最小误差为被舍弃奇异值平方之和。这也说明为何投影坐标有$r$列，而不是把样本特征向量直接当作载荷矩阵。

中心化均值和标准差仅由训练数据估计，新样本先做同样变换，再计算$z=V^{\mathsf T}(x-\bar x)$。只中心化的 PCA 偏重数值方差大的传感器，标准化后的 PCA 则赋予低方差传感器更大相对权重；二者都不是天然正确。重特征值对应的基底可以旋转，不能把单条主轴的符号或方向变化解释成物理机制变化。

\paragraph{核 PCA 的特征空间与新样本坐标。}
设正半定核$k(x,x')=\langle\varphi(x),\varphi(x')\rangle$，$K_{ij}=k(x_i,x_j)$，$K\in\R^{n\times n}$。特征空间中心化对应$K_c=HKH$，其中$H=I-\mathbf1\mathbf1^{\mathsf T}/n$。若$K_cv_j=\lambda_jv_j$、$\|v_j\|=1$且$\lambda_j>0$，则单位主轴为
\begin{equation}
 w_j=\frac1{\sqrt{\lambda_j}}\sum_i(v_j)_i
 [\varphi(x_i)-\bar\varphi],\qquad
 z_{ij}=\sqrt{\lambda_j}(v_j)_i.
\end{equation}
归一化来自$\|w_j\|^2=v_j^{\mathsf T}K_cv_j/\lambda_j=1$。对新样本$x$，先形成交叉核向量
\begin{equation}
 (k_c(x))_i=k(x_i,x)-\frac1n\sum_a k(x_a,x)
 -\frac1n\sum_bK_{ib}+\frac1{n^2}\sum_{a,b}K_{ab},
\end{equation}
再计算$z_j(x)=v_j^{\mathsf T}k_c(x)/\sqrt{\lambda_j}$。新点必须使用训练均值，不能单独重新中心化。极小特征值会放大数值误差，核宽度会改变整个几何；核空间中的好重构也不等于可以唯一还原原始传感器读数。

\paragraph{Isomap。}
Isomap 用邻域图最短路径距离$D_{ij}$近似测地距离，再执行经典多维尺度分析\cite{tenenbaum2000isomap}：
\begin{equation}
 H=I-\frac1n\mathbf1\mathbf1^{\mathsf T},\qquad
 B=-\frac12H D^{\circ2}H,\qquad Z=U_r\Lambda_r^{1/2}.
\end{equation}
$D^{\circ2}$表示逐元素平方，$U_r,\Lambda_r$取正的主特征对。算法依次构图、计算最短路径、双中心化和谱分解。图断裂时路径距离无穷，不能直接进入 Multidimensional Scaling (MDS)；跨流形捷径则会系统性缩短距离。负特征值提示估计距离不能完全对应欧氏嵌入，不应直接对其开平方。

\paragraph{双中心化为什么把距离变成内积。}
暂设$D_{ij}$确实来自中心化坐标$Z\in\R^{n\times r}$，即$Z^{\mathsf T}\mathbf1=0$。令$B=ZZ^{\mathsf T}$、$b_i=B_{ii}$，平方距离展开为
\begin{equation}
 D_{ij}^2=B_{ii}+B_{jj}-2B_{ij},\qquad
 D^{\circ2}=b\mathbf1^{\mathsf T}+\mathbf1b^{\mathsf T}-2B.
\end{equation}
由于$H\mathbf1=0$且$HBH=B$，左右乘$H$便得到$B=-HD^{\circ2}H/2$。因此谱分解不是无来由的步骤，而是在从距离反求坐标内积。若$B$半正定，保留前$r$个正特征对得到$Z=U_r\Lambda_r^{1/2}$；坐标只确定到正交旋转，中心化则固定了平移自由度。

图最短路径得到的距离未必能精确嵌入欧氏空间。截断谱解最小化的是秩约束下的 Gram 矩阵误差$\|B-ZZ^{\mathsf T}\|_F^2$，不是一般意义下的$\sum_{ij}(D_{ij}-\|z_i-z_j\|)^2$。即便图构造正确，球面等具有内在曲率的流形也不能被无失真地摊平成同维欧氏平面。Isomap 尤其依赖采样足够密、邻域不跨越折叠以及测地结构近似可展平；“存在低维流形”本身还不足以保证成功。

经典实现需要存储$n\times n$距离矩阵。Landmark 变体仅用部分代表点近似全局几何，降低计算和存储，但代表点必须覆盖各工况；选择失衡会遗漏少数状态。

\section{局部重构与谱嵌入}
\paragraph{LLE。}
局部线性嵌入先固定每个样本的邻居，求局部重构权重\cite{roweis2000lle}：
\begin{equation}
 \min_W\sum_i\left\|x_i-\sum_{j\in\mathcal N_i}W_{ij}x_j\right\|_2^2,
 \qquad \sum_jW_{ij}=1.
\end{equation}
和为一的约束使权重对整体平移不变。邻域协方差奇异时需加小的对角正则。固定$W$后，再求
\begin{equation}
 \min_Z\|(I-W)Z\|_F^2,\qquad
 Z^{\mathsf T}\mathbf1=0,\quad Z^{\mathsf T}Z=I.
\end{equation}
取$(I-W)^{\mathsf T}(I-W)$最小非平凡特征值对应向量。第一步保留局部重构关系，第二步避免坐标全部坍缩。噪声、采样密度变化或过大的邻域会破坏局部线性假设。Hessian LLE 与局部切空间对齐等变体改变局部几何估计，并非仅更换一个可视化参数\cite{donoho2003hessian,zhang2004ltsa}。

\paragraph{LLE 的两个约束优化如何求解。}
对点$i$设邻居数为$m_i$，把差向量组成$A_i=[x_j-x_i]_{j\in\mathcal N_i}\in\R^{d\times m_i}$，局部权重$w_i\in\R^{m_i}$。利用$\mathbf1^{\mathsf T}w_i=1$，重构误差化为$\|A_iw_i\|^2=w_i^{\mathsf T}C_iw_i$，其中$C_i=A_i^{\mathsf T}A_i$。拉格朗日函数取
\begin{equation}
 \mathcal J(w_i,\nu)=w_i^{\mathsf T}C_iw_i
 -2\nu(\mathbf1^{\mathsf T}w_i-1).
\end{equation}
驻点满足$C_iw_i=\nu\mathbf1$，故当$C_i$可逆时
\begin{equation}
 w_i=\frac{C_i^{-1}\mathbf1}{\mathbf1^{\mathsf T}C_i^{-1}\mathbf1}.
\end{equation}
实现应求解线性系统而非显式求逆。若邻居近似落在低维切平面，$C_i$很容易奇异，可用$C_i+\gamma I$、$\gamma>0$替代。正则化提高稳定性，但同时改变重构关系。权重允许为负，LLE 的权重矩阵不是随机游走概率矩阵。

第二步设$M=(I-W)^{\mathsf T}(I-W)\in\R^{n\times n}$。因$W\mathbf1=\mathbf1$，有$M\mathbf1=0$，不加约束时$Z=0$就是平凡最优解。约束$Z^{\mathsf T}Z=I$固定尺度、$Z^{\mathsf T}\mathbf1=0$排除常数方向；Rayleigh 商最小化给出其余最小特征值对应向量。若零空间不止一维，还需诊断连通分量或重构退化，不能机械删除一个特征向量就结束。

对新点可以在训练点中求同样的和为一重构权重，并令$z(x)=\sum_jw_j(x)z_j$。这是一种已声明的邻域插值规则，不是重新求解训练谱问题。它在已覆盖的局部区域有意义，但无法可靠外推到新工况形成的流形分支。

\paragraph{拉普拉斯特征映射。}
取对称非负邻接权$W_{ij}$，例如邻域内的高斯核权重；令$D_{ii}=\sum_jW_{ij}$、$L=D-W$。拉普拉斯特征映射最小化\cite{belkin2003laplacian}
\begin{equation}
 \frac12\sum_{i,j}W_{ij}\|z_i-z_j\|_2^2
 =\operatorname{tr}(Z^{\mathsf T}LZ),\qquad Z^{\mathsf T}DZ=I.
\end{equation}
排除常数方向后求$Lv=\lambda Dv$。图断裂会产生多个零特征值，因此非平凡方向的选择必须结合连通性。归一化拉普拉斯使用$D^{-1/2}LD^{-1/2}$，对孤立节点需单独处理，避免除零。

\paragraph{从边上的差异到广义特征值。}
本段$W\in\R^{n\times n}$是对称非负相似度，区别于上一方法允许负值的重构权重；本段的$D$是度矩阵，不是 Isomap 的距离矩阵。对单列坐标$v\in\R^n$逐项展开，有
\begin{align}
 \frac12\sum_{i,j}W_{ij}(v_i-v_j)^2
 &=\sum_iD_{ii}v_i^2-\sum_{i,j}W_{ij}v_iv_j\\
 &=v^{\mathsf T}(D-W)v\geq0.
\end{align}
在$v^{\mathsf T}Dv=1$下求驻点，$\nabla_v[v^{\mathsf T}Lv-\lambda(v^{\mathsf T}Dv-1)]=0$给出$Lv=\lambda Dv$。还需$v^{\mathsf T}D\mathbf1=0$排除常数方向。令$u=D^{1/2}v$就转成对称问题$D^{-1/2}LD^{-1/2}u=\lambda u$。低特征值表示沿强连接边缓慢变化，因而可能分开弱连接的设备群，但并不保证这些群等于故障类别。图密度、核带宽和正常工况的采样数量都可能决定这些方向。

\paragraph{扩散坐标与扩散距离。}
Diffusion Maps 将邻接权归一化为随机游走转移矩阵，以特征值的时间次幂缩放坐标\cite{coifman2006diffusion}。设图连通、各度$d_i>0$，$P=D^{-1}W$，则$P_{ij}$表示一步从$i$到$j$的概率，平稳分布为$\pi_i=d_i/\sum_ad_a$。对称性给出细致平衡$\pi_iP_{ij}=\pi_jP_{ji}$，故可选右特征函数$\psi_\ell$，满足$\sum_i\pi_i\psi_\ell(i)\psi_m(i)=\mathbf1[\ell=m]$。第一项为$\lambda_0=1,\psi_0=1$。

走$t$步后的两行概率分布之差定义扩散距离：
\begin{align}
 D_t^2(i,j)&=\sum_a\frac{[P^t_{ia}-P^t_{ja}]^2}{\pi_a}\\
 &=\sum_{\ell\geq1}\lambda_\ell^{2t}
 [\psi_\ell(i)-\psi_\ell(j)]^2.
\end{align}
第二式由$P^t_{ia}=\pi_a\sum_\ell\lambda_\ell^t\psi_\ell(i)\psi_\ell(a)$代入，并用加权正交性消去交叉项得到。因此坐标$\Psi_t(i)=(\lambda_1^t\psi_1(i),\ldots,\lambda_r^t\psi_r(i))$的欧氏距离近似扩散距离，截断误差正是未保留项之和。它衡量的是多个路径共同形成的可达性，而非一条最短路径。

时间$t$越大，$|\lambda_\ell|<1$的方向衰减越强；短时间强调局部邻域，较长时间突出缓慢变化的结构，过长则使所有点近乎重合。二部图可能出现$-1$特征值和周期振荡，可采用懒惰随机游走$(I+P)/2$消除周期性。采样密度校正改变几何与密度的相对影响，应按任务解释，不能只选最美观的嵌入。

\section{概率邻域可视化}
\paragraph{t-SNE。}
t-SNE 在原空间构造对称邻域概率$p_{ij}$，在低维空间使用重尾分布\cite{vandermaaten2008tsne}：
\begin{equation}
 q_{ij}=\frac{(1+\|z_i-z_j\|^2)^{-1}}
 {\sum_{a\ne b}(1+\|z_a-z_b\|^2)^{-1}},\qquad
 \min_Z\sum_{i\ne j}p_{ij}\log\frac{p_{ij}}{q_{ij}}.
\end{equation}
原空间高斯带宽逐点调整，使条件分布熵对应指定 perplexity。重尾低维分布缓解拥挤问题，KL 方向更重视保留高概率邻居。优化非凸，不同初始化、perplexity 和学习率可能改变结果。图中簇间空隙、面积和距离不能直接解释为真实故障严重度。

\paragraph{t-SNE 的概率、熵与梯度。}
先在高维空间定义条件邻域概率：
\begin{equation}
 p_{j\mid i}=
 \frac{\exp\bigl(-\|x_i-x_j\|^2/(2\sigma_i^2)\bigr)}
 {\sum_{a\ne i}\exp\bigl(-\|x_i-x_a\|^2/(2\sigma_i^2)\bigr)},
 \qquad j\ne i,\quad p_{i\mid i}=0.
\end{equation}
带宽$\sigma_i$按$\exp[-\sum_jp_{j\mid i}\log p_{j\mid i}]$等于给定 perplexity 调整；它描述有效邻居数而不是类别数。再令$p_{ij}=(p_{j\mid i}+p_{i\mid j})/(2n)$，便有$\sum_{i\ne j}p_{ij}=1$。

令$a_{ij}=(1+\|z_i-z_j\|^2)^{-1}$、$A=\sum_{a\ne b}a_{ab}$。去掉不含$Z$的常数，目标为
\begin{equation}
 C=-\sum_{i\ne j}p_{ij}\log a_{ij}+\log A.\label{eq:manifold-tsne-expanded-loss}
\end{equation}
式\eqref{eq:manifold-tsne-expanded-loss}右端的第一项鼓励原空间高概率邻居靠近，第二项因归一化耦合所有点而产生排斥。利用$\nabla_{z_i}\log a_{ij}=-2a_{ij}(z_i-z_j)$，并计入有序点对的两个方向，可得
\begin{equation}
 \nabla_{z_i}C=4\sum_{j\ne i}(p_{ij}-q_{ij})
 \frac{z_i-z_j}{1+\|z_i-z_j\|^2}.
\end{equation}
当$p_{ij}>q_{ij}$时梯度下降使两点靠近，反之使其分离。重尾使较远点仍能产生排斥，缓解把高维邻域挤到二维的拥挤问题，但不会恢复真实全局距离。平移、旋转不改变目标，初始化和优化路径却可能改变局部极小值；因此多次结果中稳定的邻域比某次图上的簇间距离更值得解释。

\paragraph{UMAP。}
UMAP 用局部距离尺度构造模糊邻域图，再优化低维边连接强度\cite{mcinnes2018umap}。设原图权重为$w_{ij}$，低维相似度为$q_{ij}=(1+a\|z_i-z_j\|^{2b})^{-1}$，其图交叉熵包含
\begin{equation}
 -\sum_{i<j}\left[w_{ij}\log q_{ij}+(1-w_{ij})\log(1-q_{ij})\right],
\end{equation}
其中与$Z$无关的常数省略。实际实现用边采样与负采样近似优化，不能把采样目标不加区分地当成全对全精确计算。邻居数控制局部与较大尺度结构的折中，最小距离参数影响低维点的紧密程度，不是分类间隔的直接估计。

\paragraph{UMAP 的局部尺度与吸引排斥。}
一种典型有向邻域强度在邻居上定义为
\begin{equation}
 s_{j\mid i}=\exp\left[-\frac{\max\{0,d(x_i,x_j)-\rho_i\}}{\sigma_i}\right],
\end{equation}
非邻居取零。其中$\rho_i$给出局部连接距离，$\sigma_i>0$调节有效邻域规模。以模糊并集对称化得到
\begin{equation}
 w_{ij}=s_{j\mid i}+s_{i\mid j}-s_{j\mid i}s_{i\mid j},
\end{equation}
从而即便只有一侧视对方为近邻，也可能保留该边。这与对条件概率取平均不是同一构图过程。

令$r=\|z_i-z_j\|>0$、$q=(1+ar^{2b})^{-1}$，$a,b>0$。单对损失$\ell=-w\log q-(1-w)\log(1-q)$对距离的导数为
\begin{equation}
 \frac{dq}{dr}=-\frac{2b}{r}q(1-q),\qquad
 \frac{d\ell}{dr}=\frac{2b}{r}(w-q),\qquad
 \nabla_{z_i}\ell=\frac{2b(w-q)}{r^2}(z_i-z_j).\label{eq:manifold-umap-gradient}
\end{equation}
推导式\eqref{eq:manifold-umap-gradient}只需先对$q$求导，再乘$dq/dr$。当$w>q$时下降方向拉近点对，当$w<q$时推远点对；$r=0$附近需数值稳定处理。该式描述全对全理想目标，实际边采样和负采样的频率会改变有效权重，没有对应重要性校正时不能把随机更新称为该目标的严格无偏梯度。$a,b$控制低维相似度曲线，因而紧簇大小部分来自人为设定，不是原始数据中的物理尺寸。

参数化嵌入学习显式映射，便于处理新数据；普通谱嵌入与非参数可视化通常需要邻域插值、Nystr\"om 扩展或重新拟合。全数据联合嵌入属于传导式设定，不能据此声称对未见设备具备归纳泛化能力。

\section{流形正则与工业验证}
流形正则将未标注样本几何引入监督学习\cite{belkin2006regularization}：
\begin{equation}
 \mathcal L(f)=\frac1{n_l}\sum_{i=1}^{n_l}\ell(f(x_i),y_i)
 +\lambda_A\|f\|_{\mathcal H}^2+\lambda_I\bm f^{\mathsf T}L\bm f,
\end{equation}
其中$\bm f$为全部训练侧样本的预测向量。最后一项鼓励邻近样本预测平滑；若邻域跨越真实类别边界，它会损害少数故障。类别条件图可以利用已知标签约束邻接，但不能使用测试标签构图。

\paragraph{平方损失下把函数问题化为线性系统。}
设训练侧共$n$个样本，其中$n_l$个有标签，取平方损失并假设$\lambda_A>0$。核展开$f(x)=\sum_{j=1}^na_jk(x_j,x)$足以表达最优解：任何与这些核截面正交的函数在训练点上取零，不改变经验损失和图项，却只会增加 Reproducing Kernel Hilbert Space (RKHS) 范数。因此令$a\in\R^n$、$\bm f=Ka$，$J\in\R^{n\times n}$为有标签位置为$1$的对角选择矩阵，$y$在无标签位置补零，目标化为
\begin{equation}
 \mathcal L(a)=\frac1{n_l}(Ka-y)^{\mathsf T}J(Ka-y)
 +\lambda_Aa^{\mathsf T}Ka+\lambda_Ia^{\mathsf T}KLKa.
\end{equation}
对$a$求导并除以$2$，得到
\begin{equation}
 \left(\frac1{n_l}KJK+\lambda_AK+\lambda_IKLK\right)a
 =\frac1{n_l}KJy.
\end{equation}
该系统的矩阵对称半正定。若$K$正定，$\lambda_A>0$保证唯一系数解；若$K$奇异，不能直接消去$K$，可在其非零特征子空间求解，系数虽可能不唯一，所表达的训练函数仍可相同。预测新点时使用核展开，不需要把测试标签或测试点加入训练图。

为理解图项的作用，也可暂时直接优化节点预测$f\in\R^n$。在$\frac12(f-y)^{\mathsf T}J(f-y)+\frac\gamma2f^{\mathsf T}Lf$中，驻点满足$(J+\gamma L)f=Jy$。对无标签节点$i$，有$d_if_i=\sum_jW_{ij}f_j$，即预测等于邻居的加权平均。若一个连通分量完全没有标签，其常数水平无法由此确定；若正常与故障之间有错误强边，平滑会把真实突变抹掉。几何提供的是可检验的先验，而不是无标签数据自动携带正确标签的保证。

对多工况振动数据，先统一转速、负载和通道量纲，再比较 PCA、Isomap 和图嵌入。随后改变邻居数、距离定义和随机种子，检查哪些邻近关系始终稳定。图构造和参数选择只能用训练侧数据；新设备上线时，要按事先写好的映射规则处理，不能把测试样本重新放回图中。除了邻域保持、重构误差或谱诊断，还要报告未见设备的故障召回、误报和计算成本。二维图看起来更整齐，却让异常召回下降时，它就不适合这个业务。

流形学习的算法目标和可视化边界可用下列官方材料复习；前者覆盖算法比较，后者帮助理解低维图形的解释限制。
\section*{辅助学习材料}
\begingroup
\emergencystretch=3em
\begin{itemize}
\item \textbf{Manifold learning}；作者/平台：scikit-learn Developers；英文；官方文档。重点阅读 Isomap、LLE、Spectral Embedding、t-SNE 与 MDS 小节，帮助逐一对应本章的测地距离、局部重构、图平滑和概率邻域。\href{https://scikit-learn.org/stable/modules/manifold.html}{在线阅读}。
\item \textbf{Basic UMAP Parameters}；作者/平台：UMAP-learn；英文；官方教程。比较 n\_neighbors、min\_dist、n\_components 和 metric 的示例，理解局部与全局结构、嵌入紧密度和距离选择。\href{https://umap-learn.readthedocs.io/en/latest/parameters.html}{在线阅读}。
\item \textbf{【机器学习】数据降维（Dimensionality Reduction）}；知乎；中文解说。先建立PCA基线，再比较邻域和距离保持目标；标题与主题据必应索引，原页受限。\href{https://zhuanlan.zhihu.com/p/342129669}{在线阅读}。
\item \textbf{降维算法总结（超全！附代码）}；CSDN；中文博客。核对降维的目标与新样本映射；二维可视化不能直接证明类别结构或部署泛化；标题与主题据必应索引，原页受限。\href{https://blog.csdn.net/SeafyLiang/article/details/118701759}{在线阅读}。
\item \textbf{(选修)To Learn More - Unsupervised Learning - Neighbor Embedding}；李宏毅；bilibili；中文课程。选看第93集，对照本章邻域嵌入与可视化限制。2026年10月4日官方公开元数据显示整套课程累计播放约299.5万，并非本分集播放量；已核对目录与计数，未逐帧观看。\href{https://www.bilibili.com/video/BV1Wv411h7kN/?p=93}{观看分集}。
\end{itemize}
\endgroup
\paragraph{本章小结。}
本章从同一批高维观测得到了不同的数学对象：PCA 的正交投影、核 PCA 的特征空间内积、Isomap 的测地距离、LLE 的局部重构权重、谱方法的图平滑与扩散距离，以及 t-SNE、UMAP 的邻域强度。每一种谱分解或梯度优化都服务于一个明确目标；它们不能因最终都画出二维点图而互相替代。

流形假设真正有用的地方，是把局部连续变化与有限样本下的学习联系起来；真正危险的地方，是把错误距离、采样密度或跨类别边界的邻接当成真实结构。工业验证因此要同时检查图连通性、邻域稳定性、新点映射和故障指标。低维图形提供可探索的线索，而不是类别存在性、因果关系或部署泛化的证明。

本章利用样本之间的邻近关系，下一章则利用做过多个任务的经验。例如已经在多台设备上练习过少样本诊断，能否让下一台设备只需少量标签就完成适应？流形学习要检查样本邻近是否有意义，元学习则要检查以往任务的共同规律是否适用于新任务。换了故障定义或设备机理以后，两种方法的原有假设都可能需要重审。
